\documentclass[12pt]{article}
\usepackage[utf8]{inputenc}
\usepackage[T2A]{fontenc}
\usepackage[russian]{babel}
\usepackage{amsmath}
\usepackage{array}
\usepackage{tikz}

\newcommand{\problem}{\par\noindent\textbf{Задача.} }
\newcommand{\eproblem}{\par}

\begin{document}
\problem
Рассмотрим целочисленные решения уравнения
\[
  x^2 + y^2 = 179.
\]

\begin{enumerate}
  \item Объясните, почему одно из чисел нечётно.
  \item Заполните таблицу примеров.
\end{enumerate}

\begin{center}
\begin{tabular}{|c|c|}
\hline
$x$ & $y$ \\
\hline
2 & 5 \\
\hline
\end{tabular}
\end{center}

\begin{center}
\begin{tikzpicture}[scale=0.8]
  \draw[->] (-0.2,0) -- (2.2,0);
  \draw[->] (0,-0.2) -- (0,2.2);
  \draw (0,0) -- (2,0) -- (0,2) -- cycle;
  \node at (0.75,0.75) {$179$};
\end{tikzpicture}
\end{center}
\eproblem
\end{document}
